\subsection{Simple State, Event, and History Reductions Are Bounded}

\begin{figure*}[t]
\centering
\includegraphics[width=\textwidth]{figures/4134_figure4_reduction_boundaries.png}
\caption{Reduction boundaries. The tested $(C,\dot{C},R)$ first/second moment
closure and the state-matched event-local near-pre test did not provide stable
reductions of T1. Recent-history separation was visible in some observations
but did not become a universal sign/order rule. Propagation remained outside
the current confirmatory route.}
\label{fig:reduction_boundaries}
\end{figure*}

The low-dimensional state test did not support a stable moment closure based
on $(C,\dot{C},R)$. The median incremental $R^2$ was approximately
$-6.33\times10^{-4}$ for the first moment and $-6.21\times10^{-4}$ for the
second moment, with low positive-observation fractions. These results do not
rule out all stochastic or state-dependent descriptions; they show that this
particular compact state representation did not provide a stable first- or
second-moment reduction.

The event-locality test produced a similar boundary. True transition
timestamps did not show robust near-pre excess beyond same-observation
non-event frames matched in $(C,\dot{C},R)$. The median event-minus-control
near-pre effect was approximately $-0.033$, and the positive-observation
fraction was about 0.42. This does not mean that transitions have no special
dynamics. It means that the tested state-matched near-pre aggregate did not
provide a stable event-specific explanation.

Recent path direction provided a bounded positive signal. Fourteen of 19
observations exceeded the shuffled-history median, but only 6 of 19 exceeded
the q95 null, and the sign/order of the effect was not stable across
observations. Thus, recent history remains an observation-specific boundary
rather than a universal memory or hysteresis mechanism.

